\begin{figure}[ht]
\centering
\begin{tikzpicture}[x=1cm, y=1cm, font=\small]

% Power curves for a two-sample two-sided t-test at alpha = 0.05,
% plotted against per-group sample size n on a log-x axis. Three
% standardised effect sizes (Cohen's d): 0.2, 0.5, 0.8.
% Power approximation: 1 - beta ~= Phi( d * sqrt(n/2) - z_{alpha/2} )
% with z_{0.025} = 1.96. Use the simple normal approximation; for
% n < 20 the t-correction is slight and does not change the picture.

% axes
\draw[->, thick] (0, 0) -- (9.5, 0)
  node[below right, font=\scriptsize] {sample size $n$ per group};
\draw[->, thick] (0, 0) -- (0, 5.5)
  node[left, font=\scriptsize, rotate=90, yshift=0.6cm, anchor=south]
  {power $1 - \beta$};

% x-axis: log scale n = 5 to 500
% x plot coord = 3 * (log10(n) - log10(5))
\foreach \nval/\xc in {5/0, 10/0.903, 20/1.806, 50/3.0, 100/3.903, 200/4.806, 500/6.0} {
  \draw (\xc, 0) -- (\xc, -0.1)
    node[below, font=\scriptsize] {\nval};
}
% y axis ticks
\foreach \pval/\yc in {0/0, 0.2/1, 0.4/2, 0.5/2.5, 0.6/3, 0.8/4, 1.0/5} {
  \draw (0, \yc) -- (-0.1, \yc)
    node[left, font=\scriptsize] {\pval};
}

% Reference lines: 80% power, alpha = 0.05 nominal
\draw[dashed, gray] (0, 4) -- (9.3, 4);
\node[font=\scriptsize, gray, right] at (8.5, 4.2) {power = $0.80$};

% Three power curves at d = 0.2, 0.5, 0.8. Plot domain on x is the
% log-scaled coordinate from x = 0 (n=5) to x = 6 (n=500).
% We evaluate at fine sample of x and convert: n = 5 * 10^(x/3)
% Power = 0.5 + 0.5 * erf( (d * sqrt(n/2) - 1.96) / sqrt(2) )
% (Gaussian CDF approximation). TikZ has no erf; use a hand piecewise
% polynomial approximating Phi(z) on [-3, 3], smoothed.

% d = 0.2 curve (small effect, slowly rising)
\draw[thick, engblue, smooth]
  plot[domain=0:6, samples=60]
  (\x, {5*(0.5 + 0.5*tanh(0.7886*(0.2*sqrt(5*pow(10,\x/3)/2) - 1.96)))});
\node[engblue, font=\scriptsize, left] at (5.8, 1.6) {$d = 0.2$};

% d = 0.5 curve (medium)
\draw[thick, engred, smooth]
  plot[domain=0:6, samples=60]
  (\x, {5*(0.5 + 0.5*tanh(0.7886*(0.5*sqrt(5*pow(10,\x/3)/2) - 1.96)))});
\node[engred, font=\scriptsize, right] at (3.8, 4.6) {$d = 0.5$};

% d = 0.8 curve (large, saturates quickly)
\draw[thick, engamber, smooth]
  plot[domain=0:6, samples=60]
  (\x, {5*(0.5 + 0.5*tanh(0.7886*(0.8*sqrt(5*pow(10,\x/3)/2) - 1.96)))});
\node[engamber, font=\scriptsize, right] at (1.3, 4.3) {$d = 0.8$};

\end{tikzpicture}
\caption[Power curves for a two-sample $t$-test]{Power $1 - \beta$
of a two-sample two-sided $t$-test at $\alpha = 0.05$, as a function
of per-group sample size $n$, for three standardised effect sizes
$d = \delta/\sigma$. The dashed line marks the conventional
$80\,\%$ power target. To reach it: $d = 0.8$ (large) needs $n
\approx 25$; $d = 0.5$ (medium) needs $n \approx 64$; $d = 0.2$
(small) needs $n \approx 400$. The horizontal axis is logarithmic
because $n$ scales as $1/d^{2}$. The tanh approximation to the
Gaussian CDF used here is accurate to about $0.01$ across the
plotted range and is adequate for the planning question. The plot
formalises the cost of pursuing small effects with under-powered
studies, which is the statistical engine behind the winner's curse
named in section 8.6.}
\label{fig:vol01:ch08:power-curve}
\end{figure}
